\documentclass[11pt]{article}

% basic packages
\usepackage[margin=1in]{geometry}
\usepackage[pdftex]{graphicx}
\usepackage{amsmath,amssymb,amsthm}
\usepackage{custom}
\usepackage{lipsum}
\usepackage{booktabs}
\usepackage{enumitem}

% page formatting
\usepackage{fancyhdr}
\pagestyle{fancy}

\renewcommand{\sectionmark}[1]{\markright{\textsf{\arabic{section}. #1}}}
\renewcommand{\subsectionmark}[1]{}
\lhead{\textbf{\thepage} \ \ \nouppercase{\rightmark}}
\chead{}
\rhead{}
\lfoot{}
\cfoot{}
\rfoot{}
\setlength{\headheight}{14pt}

\linespread{1.03} % give a little extra room
\setlength{\parindent}{0.2in} % reduce paragraph indent a bit
\setcounter{secnumdepth}{2} % no numbered subsubsections
\setcounter{tocdepth}{2} % no subsubsections in ToC

\begin{document}

% make table of contents
\newpage
\microtoc
\newpage

% main content
\section{The Schr\"odinger Equation}
\subsection{Wave Functions}
\textbf{Example.} An extremely common continuous distribution is the Gaussian distribution:
\(\rho(x) = Ae^{-\lambda (x-a)^{2}}\). 
\\

\noindent To normalize the distribution, we demand: 
\[
    \int ^{\infty}_{-\infty} Ae^{-\lambda (x-a) ^{2}}dx = 1
\]
Translating to be centered around the origin does not change the integral. Then let us
substitute \(\sqrt{\lambda}x = y\) we have
\[
    \frac{\sqrt{\lambda}}{A} = \int^{\infty}_{-\infty} e^{-y^2}dy = \sqrt{\pi}
\]
So \(A = \sqrt{\frac{\lambda}{\pi}}\). Meanwhile, to find the expected value, make the
substitution \((x-a) = y \), giving us 
\[
    \langle x \rangle = A\int^{\infty}_{-\infty}ye^{-\lambda y^2}dy 
    + aA\int^{\infty}_{-\infty}e^{-\lambda y^2}dy =  -I_1 + I_1 + a = a   
\]
So we can expect the particle to be located at the middle of the gaussian distribution. Meanwhile, 
\[
    \langle x^2 \rangle = A\int^{\infty}_{-\infty}y^2e^{-\lambda y^2}dy 
    + A\int^{\infty}_{-\infty}2ay e^{-\lambda y^2}dy + A\int^{\infty}_{-\infty}a^2 e^{-\lambda y^2}dy = 
    \frac{1}{2\sqrt{\lambda}} + a^2
\]

\noindent Therefore, the standard deviation \(\sigma_x \) is \(\frac{1}{2\sqrt{\lambda}}\)
\\



\noindent
\textbf{Example.} Here, we'll examine the time dependence of \(\langle p \rangle\). 
\\

\noindent The derivative is:
\[
    \frac{d\langle p \rangle }{dt} = -i\hbar \int \frac{\partial}{\partial t}\left(\Psi^* \frac{\partial \Psi}{\partial x}\right)dx
\]
using Schr\"odinger's equation along with the product rule. Expanding using the product rule and substituting in 
Schr\"odinger's equation gives us 
\[
    \int \frac{\hbar^2}{2m}\left(-\frac{\partial^2 \Psi}{\partial x^2}\frac{\partial \Psi}{\partial x} 
    + \Psi^* \frac{\partial^3 \Psi}{\partial x^3} \right) + V\Psi^* \frac{\partial \Psi}{\partial x} 
    - \Psi^*\frac{\partial (V\Psi)}{\partial x}dx
\]

\noindent Integrating by parts twice on the left term and applying the boundary conditions that \(\frac{\partial\Psi}{\partial x}\) and 
\(\frac{\partial \Psi^*}{\partial x}\) go to 0 at infinity cancels the two out. Expanding the right side, we have
\[
    \frac{\partial (V\Psi)}{\partial x} = \Psi \frac{\partial V}{\partial x} + V \frac{\partial \Psi}{\partial x} 
\]
so we are left with 
\[
    \frac{d\langle p \rangle }{dt} = \int \Psi^* \left(- \frac{\partial V}{\partial x}\right)\Psi dx = \langle -\frac{\partial V}{\partial x}\rangle
\]
\noindent This is Ehrenfest's theorem for Newton's second law. 
\\

\noindent
\textbf{Griffiths Problem 1.11 and 1.12} A classical approach to creating a wave function is by directly modeling the probability distribution
of a classical particle using kinematics. Suppose a particle of mass \(m\) is in a classical harmonic oscillator of potential \(V(x) = \frac{kx^2}{2}\) and
has energy \(E\). Let the fraction of time the particle spends at position \(x\) be 
\[
    \rho(x)dx = \frac{dt}{T} = \frac{(dt/dx)dx}{T} = \frac{1}{v(x)T}dx
\]
where T is the total period. Then the quantity \(\rho(x) = \frac{1}{v(x)T}\) is a close analogy to \(|\Psi|^2\). (Careful, \(T\) is the time it takes to
travel from one side to the other, so it is half of the period)
\\

\noindent Using conservation of energy, \(v(x) = \sqrt{\frac{2(E-\frac{kx^2}{2})}{m}}\). Then we have \(\rho(x) = \frac{1}{\pi\sqrt{(a^2-x^2)}}\). Integrating from
a to -a, \(\int^a_{-a} \frac{dx}{\pi\sqrt{a^2-x^2}} = \frac{1}{\pi}\left[\arcsin(\frac{x}{a})\right]^a_{-a} = 1\) so this "wave function" satisfies the normalization
constraint. Meanwhile we have
\[
    \langle x\rangle = \int^a_{-a}\frac{x}{\pi \sqrt{a^2-x^2}}dx = \left[\frac{-\sqrt{a^2-x^2}}{\pi}\right]^a_{-a} = 0
\]

\[
    \langle x^2\rangle = \int^a_{-a}\frac{x^2}{\pi \sqrt{a^2-x^2}}dx = \frac{a^2}{\pi}\int^{\frac{\pi}{2}}_{-\frac{\pi}{2}} \sin^2{\left(\theta\right)}d\theta = \frac{a^2}{2}
\]

\noindent Thus, the standard deviation \(\sigma_x\) is \(\frac{a}{\sqrt{2}}\). Meanwhile, if we were to do the same with momentum, we can write
\[
    \rho(p)dp = \frac{dt}{T} = \frac{(dt/dp)dp}{T} = \frac{dp}{T|\dot{p}|}
\]
From conservation of energy, we have 
\[
    \dot{p} = -\frac{dV}{dx} = -kx = \sqrt{2Ek-\omega^2p^2}
\]
Then the "momentum wave" is 
\[
    \rho(p) = \frac{1}{(\pi/\omega)\sqrt{2kE - w^2p^2}} = \frac{1}{\pi \sqrt{2mE - p^2}} = \frac{1}{\pi \sqrt{p_0^2 - p^2}}
\]

\noindent This is in the exact same form as the position wave, so we already know 
\[
    \langle p \rangle = 0
\]
\[
    \langle p^2 \rangle = \frac{p_0^2}{2}
\]

\noindent Then using the standard deviations, we can see that 
\[
    \sigma_x \sigma_p = \frac{p_0x_0}{2}
\]

\noindent This should be setting off alarm bells! From the Heisenberg Uncertainty Principle, we must have
\[
    \sigma_x \sigma_p \geq \frac{\hbar}{2}
\]
but with our classical "wave," we can pick any arbitrarily small starting position. Clearly, this classical representation does not work. 

\subsection{Time Independent Schr\"odinger Equation}
\textbf{Example.} Suppose the separation constant \(E\) is allowed to have an imaginary component. While an imaginary energy is physically
impossible, there isn't immediately anything mathematically preventing this. 
\\
\\
Let's consider a wavefunction whose energy has an imaginary 
component \(E = E_0 + i\Gamma\) where \(\Gamma\) is strictly real. Then the wave function is \(\Psi(x,t) = \psi(x)e^{-i\hbar(E_0 + i\Gamma)/t}\).
The normalization condition then gives us

\[
    1 = \int \Psi(x,t) = \int \psi^*\psi e^{2\Gamma t / \hbar} dx = e^{2\Gamma t / \hbar}
\]

\noindent But in order for this to hold for all time \(t\), this implies \(\Gamma = 0\). Thus, as expected physically, the total energy must be 
strictly real.
\\

\noindent \textbf{Example.} Suppose we have an imaginary solution \(\psi(x)\). Then conjugating both sides of the time independent Schr\"odinger equation give us 
\[
    -\frac{\hbar^2}{2m}\frac{d(\psi^{*})^2}{dx^2} + V\psi^* = E\psi^*
\]
since E (and consequently, V) must be strictly real. So the conjugate, real wave  \(\psi^*\)  is also a solution. But any linear combination of solutions
is also a solution so both 
\[
  \psi_1 = (\psi + \psi^*)  
\]
\[
    \psi_2 = i(\psi - \psi^*)
\]
are also solutions. Then we can rewrite \(\psi = \frac{1}{2}\psi_1 + \frac{1}{2i}\psi_2\). What this means is any time we have an imaginary solution to the 
Schr\"odinger solution, we can simply work with its real complex conjugate, and then if needed, convert back to the imaginary solution.
\\

\noindent \textbf{Example.} If \(V(x)\) is an even function we can see that if \(\psi(x)\) is a solution, then
\[
    -\frac{\hbar^2}{2m}\frac{d\psi(-x)^2}{dx^2} + V(-x)\psi(-x) = E\psi^*
\]
\[
    -\frac{\hbar^2}{2m}\frac{d\psi(-x)^2}{dx^2} + V(x)\psi(-x) = E\psi^*
\]

\noindent \(\psi(-x)\) is also a solution. But the linear combinations \(\psi(x) + \psi(-x)\) and \(\psi(x) - \psi(-x)\) are also solutions. \(\psi(x) + \psi(-x)\) is 
even and \(\psi(x) - \psi(-x)\). So in this case we can always just choose to work with even or odd wavefunctions.

\subsection{Infinite Potential Well}
\textbf{Example.} We'll examine the stationary states of the infinite potential well. 
\\
The spatial component is \( \psi_{n}(x) = \sqrt{\frac{2}{a}}\sin{\left( \frac{n\pi}{a}x \right)} \) so we have 

\[ 
    \langle x \rangle = \frac{a}{2}
\]


\[ 
    \langle x^{2} \rangle = a^{2}\left( \frac{1}{3} - \frac{1}{2n^{2}\pi^{2}}\right)
\]

\noindent We can easily compute \( \langle p \rangle \) and \( \langle p^{2} \rangle \)using Ehrenfest's theorem. Using \( \langle p \rangle = m \frac{d\langle x\rangle}{dt} \) and
\( E_{n} = \frac{\langle p^{2} \rangle}{2m} \) then we have


\[ 
    \langle p \rangle = 0 
\]


\[ 
    \langle p^{2} \rangle  = \frac{n^{2} \pi^{2} \hbar^2}{a^{2}}
\]

\noindent This lines up with the uncertainty principle, as we have

\[ 
    \sigma_{x} \sigma_{p} =  \sqrt{ \left(\frac{a^{2}}{12} - \frac{a^{2}}{2n^{2}\pi^{2}} \right) } \frac{n\pi\hbar}{a} = \hbar\sqrt{\frac{n^{2}\pi^{2}}{2} - \frac{1}{2}} 
\]

\noindent which is minimized at the ground state \( \sigma_{x} \sigma_{p} = \hbar \sqrt{\frac{\pi^{2}}{2} - \frac{1}{2}} \approx 0.5678\hbar \geq \frac{\hbar}{2} \).
\\

\noindent \textbf{Example.} In principle, we can have any wave function as our starting condition as long as it's normalizable.
Consider the wave function 
\[
    \Psi(x,0) = 
    \begin{cases}
        Ax & 0 \leq x \leq \frac{a}{2} \\
        A(a - x) & \frac{a}{2} \leq x \leq a 
    \end{cases}  
\]
\\
In order to be normalizable, we need \(\int \Psi^2 dx = 2(\frac{A^2a^3}{24}) = 1\) which means \(A = \frac{2 \sqrt{3} }{ a^{3/2}}\). Note that this function is even about
the center of the well, so only the odd \(n\) components give us anything. Then using Fourier's trick: 
\[
    c_n = \int^{\frac{a}{2}}_0 Ax\sqrt{\frac{2}{a}}\sin{\left(\frac{n\pi}{a}x\right)} = \frac{4 \sqrt{6}}{n^2 \pi^2}\sin{\left(\frac{n\pi}{2}\right)}
\]

\noindent So the wavefunction as a function of time is simply
\[
    \Psi(x,t) = \frac{4 \sqrt{6}}{\pi^2}\sum_{n=1}^{\infty} \frac{(-1)^{(n-1)/2}}{n^2}\sin{\left(\frac{n\pi x}{a}\right) e^ {-iE_n t / \hbar}} 
\]

\noindent The probability of finding energy \(E_1\) is \(\frac{96}{\pi^4} \approx 0.9855\) and the expectation of energy is 
\(\langle H \rangle = \frac{96 E_1}{\pi^4} \sum_{n=1}^{\infty}\frac{1}{n^2}  = \frac{12}{\pi^2}E_1 \approx 1.216E_1 \). Within the wave function, the \(\psi_1\) terms dominate
and this is reflected in the expectation of energy. 


\end{document}